\chapter{数字集成电路后端设计流程与全局布线技术基础}
\label{cha:backend-flow-and-global-routing-fundamentals}

\section{数字集成电路后端设计全流程}

数字集成电路后端设计（Physical Design）是连接前端逻辑设计与芯片制造的核心环节，其核心目标是将前端综合生成的门级网表，转换为满足时序、功耗、面积与可制造性约束的物理版图，最终输出可用于流片的图形数据库文件（Graphic Database System II, GDSII）。随着集成电路工艺进入7\,nm及以下先进节点，芯片规模达到百亿级晶体管级别，后端设计复杂度显著提升。全局布线（Global Routing）作为后端设计中影响时序收敛、绕通率与良率的关键步骤，其技术优化已成为后端设计的重要研究方向。

完整的数字集成电路后端设计全流程可分为7个核心阶段。各阶段前后衔接、迭代优化，共同保证芯片设计的最终交付。图~\ref{fig:chap2-flow-overview-placeholder}给出了流程框图占位。

\begin{figure}[!htbp]
    \centering
    \fbox{\parbox{0.88\textwidth}{
            \centering
            图位占位：图2-1 数字集成电路后端设计全流程框图\\
            要求：按顺序标注各核心阶段与输入输出数据流向。
        }}
    \caption{数字集成电路后端设计全流程框图（占位）}
    \label{fig:chap2-flow-overview-placeholder}
\end{figure}

\subsection{设计导入与数据准备}

设计导入是后端设计的起始阶段，核心是完成全流程所需设计数据与工艺数据的校验与导入，避免前端设计错误流入后续环节并造成高额迭代成本。该阶段核心输入包括：
\begin{enumerate}
    \item \textbf{门级网表}：由前端逻辑综合生成，包含标准单元、宏单元（Static Random Access Memory, SRAM；Intellectual Property, IP核等）的实例化与连接关系，需经过形式验证以确保与寄存器传输级（Register Transfer Level, RTL）设计功能一致；
    \item \textbf{时序约束文件}：通常为Synopsys设计约束（Synopsys Design Constraints, SDC）格式，定义芯片工作频率、输入输出延迟、时钟特性与时序例外，是全流程时序优化的核心依据；
    \item \textbf{工艺库文件}：包含代工厂提供的时序库（\texttt{.lib}）、物理库（Library Exchange Format/Design Exchange Format, LEF/DEF）及设计规则检查（Design Rule Check, DRC）/版图与原理图一致性检查（Layout Versus Schematic, LVS）规则文件，定义标准单元与宏单元的物理尺寸、引脚位置、时序特性、功耗参数及制造规则约束；
    \item \textbf{物理约束文件}：定义芯片核心区尺寸、输入输出（Input/Output, IO）单元位置、宏单元摆放约束与电源域划分等物理实现要求。
\end{enumerate}

该阶段的核心输出是经过一致性与完整性校验的设计数据库，为后续布图规划提供可靠基础。

\subsection{布图规划（Floorplan）}

布图规划是后端设计中决定芯片整体性能的关键环节，其目标是完成芯片整体物理布局规划，确定宏单元、IO单元与电源网络位置，为后续标准单元布局与布线预留资源，并在面积、时序与功耗之间实现折中。

布图规划的核心内容包括：
\begin{enumerate}
    \item \textbf{宏单元布局}：针对SRAM、锁相环（Phase Locked Loop, PLL）、高速接口IP等硬核宏单元，根据接口方向、时序需求与功耗特性进行摆放。实践中通常将宏单元沿核心区边缘排布，以减少内部通道阻挡并缩短关键互连长度；
    \item \textbf{IO单元与电源环设计}：在芯片边界排布IO与静电放电（Electrostatic Discharge, ESD）防护单元，同时构建电源环（Power Ring）以保障核心区供电稳定；
    \item \textbf{电源地网络设计}：在核心区构建横竖交错的电源地网格（Power Delivery Network, PDN / Power Grid），依据功耗分布调整线宽与密度，保证压降（IR Drop）与电迁移（Electromigration, EM）满足约束；
    \item \textbf{布局区域划分}：依据电源域与功能模块划分核心区，定义各区域标准单元利用率上限，为后续布局与布线保留资源裕量。
\end{enumerate}

图~\ref{fig:chap2-floorplan-placeholder}给出了典型布图规划示意占位。

\begin{figure}[!htbp]
    \centering
    \fbox{\parbox{0.88\textwidth}{
            \centering
            图位占位：图2-2 典型数字芯片布图规划示意图\\
            要求：标注宏单元、IO环、电源网格、核心布局区与布线通道。
        }}
    \caption{典型数字芯片布图规划示意图（占位）}
    \label{fig:chap2-floorplan-placeholder}
\end{figure}

布图规划质量直接影响后续布局布线收敛性。不合理的布图规划会诱发严重拥塞、时序违例与电源噪声问题，甚至导致回退重做，显著拉长设计周期。

\subsection{布局（Placement）}

布局环节的核心目标是将标准单元摆放到布图规划定义的核心区合法位置，在满足不重叠与可制造约束的前提下，最小化总线长、优化时序并均衡拥塞，为后续时钟树综合与布线奠定基础。

标准布局流程通常分为三个阶段：
\begin{enumerate}
    \item \textbf{全局布局（Global Placement）}：不强制满足物理合法性，以线长与时序目标为主进行初始分布，常采用非线性规划或二次规划；
    \item \textbf{合法化（Legalization）}：将单元对齐到标准单元行并消除重叠，尽量减小位移以保留全局布局收益；
    \item \textbf{详细布局（Detailed Placement）}：通过局部交换、平移和尺寸调整等操作，进一步优化局部线长、时序与拥塞。
\end{enumerate}

\subsection{时钟树综合（Clock Tree Synthesis, CTS）}

时钟网络负载大、翻转频繁，其传输质量直接影响芯片最高工作频率与时序收敛。时钟树综合（Clock Tree Synthesis, CTS）的核心目标是构建从时钟源到各时序单元时钟引脚的时钟分发网络，在满足时钟偏移（Skew）和时钟延迟（Latency）约束的同时，控制功耗与面积并提升鲁棒性。

CTS核心流程包括：
\begin{enumerate}
    \item \textbf{拓扑规划}：依据时钟域划分与触发器分布选择拓扑结构，如H树、X树、鱼骨型与平衡缓冲树；
    \item \textbf{缓冲器插入与树构建}：通过插入缓冲器分割长线并平衡各路径延迟，常见方法包括延迟合并嵌入（Deferred Merge Embedding, DME）与Prim-Dijkstra折中策略；
    \item \textbf{时钟树优化}：利用有用偏斜（Useful Skew）改善建立/保持违例，并通过门控时钟单元优化降低动态功耗；
    \item \textbf{时钟树签核}：提取寄生参数并完成静态时序分析，验证偏移、延迟与脉宽等指标。
\end{enumerate}

CTS完成后，主路径时序结构基本固定，后续布线阶段可调整空间有限，因此CTS是后端时序收敛的重要里程碑。

\subsection{布线（Routing）}

布线是后端设计中耗时最长、复杂度最高的环节。其目标是在金属层上完成信号网与电源网的物理走线，并满足DRC、时序、信号完整性与可制造性要求。

工业界布线流程通常包括两个阶段：
\begin{enumerate}
    \item \textbf{全局布线（Global Routing）}：将核心区划分为全局网格单元（Global Cell, G-Cell），在不确定具体几何细节前提下规划线网大致路径，估计拥塞、线长与延迟，并完成初步层分配；
    \item \textbf{详细布线（Detailed Routing）}：在全局路径指导下确定精确走线位置、线宽与过孔，严格满足工艺规则并进一步优化串扰、时序与功耗。
\end{enumerate}

图~\ref{fig:chap2-metal-stack-placeholder}给出了金属层堆叠示意占位。

\begin{figure}[!htbp]
    \centering
    \fbox{\parbox{0.88\textwidth}{
            \centering
            图位占位：图2-5 芯片金属层堆叠结构示意图\\
            要求：展示上下层金属的水平/垂直交替方向、过孔结构、底层与高层金属线宽差异。
        }}
    \caption{芯片金属层堆叠结构示意图（占位）}
    \label{fig:chap2-metal-stack-placeholder}
\end{figure}

全局布线的路径质量直接决定详细布线的绕通率、时序收敛与迭代次数，是本章后续讨论重点。

\subsection{物理验证与签核（Signoff）}

布线完成后进入签核阶段，目标是验证物理版图满足全部设计与制造约束，确保流片后功能正确、性能达标且良率可控。核心内容包括：
\begin{enumerate}
    \item \textbf{DRC检查}：验证线宽、线距、过孔尺寸与包围区等制造规则；
    \item \textbf{LVS检查}：验证版图连接关系与前端网表一致性；
    \item \textbf{时序签核}：提取寄生并执行静态时序分析（Static Timing Analysis, STA），覆盖多工艺角与工作环境；
    \item \textbf{功耗与电源完整性签核}：评估动态/静态功耗并验证IR Drop与EM约束；
    \item \textbf{可制造性设计检查}：执行可制造性设计（Design for Manufacturability, DFM）优化，如冗余过孔与光刻友好修正。
\end{enumerate}

\subsection{先进工艺下后端设计的核心挑战}

进入7\,nm及以下节点后，后端设计对全局布线提出更高要求，主要挑战包括：
\begin{enumerate}
    \item \textbf{布线资源紧张与拥塞加剧}：晶体管密度持续提升而布线资源增长受限，拥塞问题更加突出；
    \item \textbf{时序收敛难度上升}：金属电阻增大导致长线RC延迟显著增加，关键路径优化压力上升；
    \item \textbf{信号完整性与电源完整性风险上升}：串扰、IR Drop与EM问题更敏感，需在全局阶段提前感知；
    \item \textbf{设计规模与复杂度快速增长}：超大规模线网使串行算法效率不足，并行化与加速计算成为必然。
\end{enumerate}

\section{全局布线技术原理与标准流程}

全局布线是详细布线前的关键预布线环节。其核心是在不确定具体几何形态的前提下，为全部线网生成可行全局路径并预分配资源，在兼顾线长、时序与拥塞的同时，为详细布线提供高质量指导，降低后续迭代成本。

\subsection{全局布线的核心定义与输入输出}

\subsubsection{全局布线的网格模型}

工业界主流全局布线工具普遍采用G-Cell网格模型：将核心区划分为等大小矩形网格单元。每个G-Cell边界对应布线通道，通道具有固定容量（Capacity），由金属层可用资源、线宽与线距等工艺参数共同决定。当实际需求超过容量时即产生拥塞（Congestion）。

图~\ref{fig:chap2-gcell-model-placeholder}给出了网格模型占位。

\begin{figure}[!htbp]
    \centering
    \fbox{\parbox{0.88\textwidth}{
            \centering
            图位占位：图2-6 全局布线G-Cell网格模型示意图\\
            要求：标注G-Cell边界、布线容量、已占用容量、拥塞区域定义与可视化方式。
        }}
    \caption{全局布线G-Cell网格模型示意图（占位）}
    \label{fig:chap2-gcell-model-placeholder}
\end{figure}

\subsubsection{全局布线的核心目标}

全局布线是多目标优化问题，主要目标包括：
\begin{enumerate}
    \item \textbf{绕通率优先}：保证线网可行路径并控制整体拥塞；
    \item \textbf{时序优化}：优先保障关键路径线网延迟；
    \item \textbf{线长最小化}：在满足约束前提下降低总线长；
    \item \textbf{资源均衡}：均衡不同区域与不同层资源占用；
    \item \textbf{可制造性优化}：控制过孔数量并合理使用高层长距离走线。
\end{enumerate}

\subsubsection{全局布线的输入与输出}

全局布线输入通常包括：布局后的DEF、网表与SDC、工艺LEF，以及布图规划与PDN障碍信息。输出通常包括：
\begin{itemize}
    \item 线网的G-Cell级路径序列；
    \item 全芯片拥塞分布图；
    \item 金属层分配结果；
    \item 线长与延迟预估结果。
\end{itemize}

工业界流程通常分为树生成（Tree Generation）、模式布线与层分配（Pattern Routing and Layer Assignment）、迷宫布线（Maze Routing）三个阶段。

\subsection{树生成阶段（Tree Generation）}

树生成阶段将多端线网转换为连通树拓扑，确定引脚连接关系与Steiner点（Steiner Point）位置，为后续路径规划提供拓扑基础。在曼哈顿平面场景下，多端线网线长最优解通常对应直角Steiner最小树（Rectilinear Steiner Minimum Tree, RSMT）。

\subsubsection{Flute算法：主流RSMT构建方法}

Flute（Fast Lookup Table Based Rectilinear Steiner Minimal Tree）算法是工业界与学术界应用广泛的RSMT构建算法，在主流EDA流程中有良好工程实践。其核心优势是对小引脚数线网可高效获取高质量拓扑，并对大引脚数线网通过分治策略保持较优质量与速度平衡。

Flute流程包括：
\begin{enumerate}
    \item \textbf{预处理阶段}：针对小规模引脚组合建立查找表（Look-Up Table, LUT）；
    \item \textbf{在线处理阶段}：按引脚规模选择查表或分治合并策略，输出近似最优RSMT。
\end{enumerate}

此外，Flute支持增量更新，适用于工程变更指令（Engineering Change Order, ECO）场景下的局部树重构。

\subsubsection{时序驱动的树生成算法}

线长最优并不必然等价于时序最优。对关键路径而言，需要在树生成阶段引入时序驱动目标。主流方法包括：
\begin{enumerate}
    \item \textbf{Prim-Dijkstra折中算法}：通过权重系数在最小生成树（Minimum Spanning Tree, MST）与最短路径树（Shortest Path Tree, SPT）之间折中；
    \item \textbf{缓冲器插入感知树生成}：在拓扑构建时引入缓冲器位置与代价；
    \item \textbf{时钟树专用方法}：如DME算法，重点优化偏移一致性。
\end{enumerate}

\subsection{Pattern Routing与层分配（Layer Assignment）}

Pattern Routing通过预定义模式将树边快速映射到G-Cell网格，是全局布线效率的核心来源。其特点是无需全图搜索即可覆盖大部分线网，再将少量困难线网交由Maze Routing处理。

\subsubsection{常用Pattern Routing走线模式}

常见模式包括：
\begin{enumerate}
    \item \textbf{L型模式}：单拐点、实现简单，适合低拥塞短连线；
    \item \textbf{Z型（N型）模式}：双拐点，保持曼哈顿最短线长且具备更强绕障能力；
    \item \textbf{十字型模式}：面向Steiner点的多分支连接；
    \item \textbf{X型模式}：面向支持45度走线的工艺场景；
    \item \textbf{鱼骨型模式}：面向总线类并行线网的协同规划。
\end{enumerate}

图~\ref{fig:chap2-patterns-placeholder}给出了模式示意占位。

\begin{figure}[!htbp]
    \centering
    \fbox{\parbox{0.88\textwidth}{
            \centering
            图位占位：图2-9 常用Pattern Routing走线模式示意图\\
            要求：包含L型、Z型、十字型、X型、鱼骨型Pattern及适用场景对比。
        }}
    \caption{常用Pattern Routing走线模式示意图（占位）}
    \label{fig:chap2-patterns-placeholder}
\end{figure}

\subsubsection{Pattern Routing执行流程}

Pattern Routing通常遵循时序优先原则：
\begin{enumerate}
    \item 按时序优先级对线网排序；
    \item 逐网逐边选择合适Pattern生成初始路径；
    \item 校验容量并在必要时更换Pattern以规避拥塞；
    \item 完成后进行局部后处理以均衡拥塞并压缩线长。
\end{enumerate}

\subsubsection{层分配（Layer Assignment）}

层分配将全局路径映射到具体金属层，需要在时序、拥塞与过孔成本间折中。一般而言，低层金属更适合局部短连线，高层金属更适合长距离关键走线。

层分配核心目标包括：
\begin{enumerate}
    \item 最小化总延迟并优先保障关键线网；
    \item 均衡各层拥塞；
    \item 控制过孔数量；
    \item 满足DRC约束。
\end{enumerate}

\paragraph{基于动态规划的层分配算法}

动态规划方法可用于单线网层分配最优化。设路径分段为$N$，金属层数为$M$，定义状态$dp[i][k]$表示第$i$段分配到第$k$层时的最小代价，则状态转移为：
\begin{equation}
    dp[i][k] = \min_{1 \leq j \leq M}\left( dp[i-1][j] + cost_{switch}(j,k) + cost_{layer}(i,k) \right)
\end{equation}
其中，$cost_{switch}$表示层切换代价，$cost_{layer}$表示当前层延迟与拥塞综合代价。该方法复杂度为$O(N\cdot M^2)$，在工程可接受范围内且效果稳定。

\begin{table}[!htbp]
    \centering
    \caption{动态规划层分配算法伪代码（占位）}
    \label{tab:chap2-dp-layer-assignment-placeholder}
    \fbox{\parbox{0.88\textwidth}{
            表位占位：后续补充动态规划层分配算法伪代码。
        }}
\end{table}

\subsection{Maze Routing阶段}

Maze Routing用于处理Pattern Routing未完成线网，通过网格搜索提升最终绕通率。其优势是可行性强，但直接全局搜索代价较高，因此工程中通常采用分层或局部化策略。

\subsubsection{经典Lee迷宫布线算法}

Lee算法本质是广度优先搜索（Breadth-First Search, BFS），在可行时保证最短路径。其步骤包括网格初始化、波前扩展与路径回溯。尽管鲁棒性强，但在大规模网格上计算与存储开销较高。

\begin{figure}[!htbp]
    \centering
    \fbox{\parbox{0.88\textwidth}{
            \centering
            图位占位：迷宫布线水波法示意图。
        }}
    \caption{迷宫布线水波法示意图（占位）}
    \label{fig:chap2-lee-wave-placeholder}
\end{figure}

\subsubsection{分层迷宫布线：粗粒度与细粒度折中方案}

分层迷宫布线将求解拆分为粗粒度全局搜索与细粒度局部搜索两阶段：
\begin{enumerate}
    \item \textbf{粗粒度阶段}：通过超级网格（Super Cell）降低状态规模并快速规划全局通路；
    \item \textbf{细粒度阶段}：在局部窗口内完成精细路径搜索，失败时扩大窗口或回退重规划。
\end{enumerate}

该方案可显著降低搜索规模，在效率与质量之间取得较好平衡。

\begin{figure}[!htbp]
    \centering
    \fbox{\parbox{0.88\textwidth}{
            \centering
            图位占位：图2-12 分层迷宫布线算法示意图\\
            要求：展示粗粒度超级网格路径规划与细粒度局部网格搜索对比。
        }}
    \caption{分层迷宫布线算法示意图（占位）}
    \label{fig:chap2-hier-maze-placeholder}
\end{figure}

\subsubsection{迷宫布线优化技术}

常见优化包括：
\begin{enumerate}
    \item 双向搜索（Bidirectional Search）；
    \item 时序驱动代价函数；
    \item 拥塞驱动代价函数。
\end{enumerate}

\section{全局布线的优化变种算法}

随着先进工艺下设计规模持续增长，传统流程在效率与质量上都面临瓶颈。工业界与学术界围绕搜索策略、并行调度与目标函数构建提出了多类优化变种。

\subsection{A*算法与启发式搜索优化}

A*算法是在Lee算法基础上引入启发式函数的有指导搜索方法，能够在保证最优性的前提下显著减少搜索节点数量，是当前迷宫布线中的重要技术。

\subsubsection{A*算法核心原理}

A*的评价函数定义为：
\begin{equation}
    f(n) = g(n) + h(n)
\end{equation}
其中，$g(n)$为起点到当前节点的实际代价，$h(n)$为到终点的启发式估计。为保证全局最优，$h(n)$应满足可采纳性。对曼哈顿布线，常用启发式函数为：
\begin{equation}
    h(n) = |x_n - x_{end}| + |y_n - y_{end}|
\end{equation}

\subsubsection{A*算法实现流程}

A*标准流程包括：
\begin{enumerate}
    \item 初始化开放集合（Open）与关闭集合（Closed）；
    \item 每轮取$f(n)$最小节点扩展；
    \item 更新邻接节点代价与父指针；
    \item 抵达终点后回溯得到路径。
\end{enumerate}

图~\ref{fig:chap2-astar-vs-lee-placeholder}给出了A*与Lee搜索范围对比占位。

\begin{figure}[!htbp]
    \centering
    \fbox{\parbox{0.88\textwidth}{
            \centering
            图位占位：图2-13 A*算法与Lee算法搜索范围对比示意图\\
            要求：体现启发式引导带来的搜索范围缩小效果。
        }}
    \caption{A*算法与Lee算法搜索范围对比示意图（占位）}
    \label{fig:chap2-astar-vs-lee-placeholder}
\end{figure}

\subsection{全局布线并行化算法}

面向千万级以上线网规模，串行全局布线已难满足工程时限。并行化全局布线的核心挑战在于资源依赖与冲突控制：不同线网对共享G-Cell容量的竞争会影响拥塞估计与最终绕通质量。

因此，并行算法通常需要同时设计：
\begin{enumerate}
    \item \textbf{任务划分机制}：按区域、按线网簇或按优先级分层划分任务；
    \item \textbf{冲突消解机制}：采用容量预留、回滚重路由或代价重加权；
    \item \textbf{同步与收敛机制}：通过批次提交与全局拥塞反馈实现稳定收敛。
\end{enumerate}

\section{本章小结}

本章系统介绍了数字集成电路后端设计全流程与全局布线关键技术。首先从流程层面梳理了设计导入、布图规划、布局、CTS、布线与签核等关键阶段，明确了全局布线在后端流程中的枢纽作用；随后从算法层面分析了树生成、Pattern Routing与层分配、Maze Routing等标准流程，并进一步讨论了A*启发式搜索与并行化优化方向。

总体而言，先进工艺下全局布线面临资源紧张、时序收敛困难与规模爆发等挑战。围绕拥塞感知、时序驱动与并行加速的系统化优化，将是后续章节展开算法研究与工程实现的基础。
